\paragraph{R4-O1, R4-W1 (Core contribution feels like a re-abstraction of constrained retrieval).}

\begin{itemize}[leftmargin=1.2em, itemsep=1pt, topsep=2pt]
\item \textit{O1: ``\dots currently it feels more like a re-abstraction of constrained retrieval. The authors are encouraged to further clarify this point.''}
\item \textit{W1: ``The distinction between counterfactual querying and existing constrained retrieval methods is still not fully clear.''}
\end{itemize}
\smallskip

\noindent\textbf{Response.}
\hw{We agree the distinction from constrained retrieval should be explained more explicitly, and we do so below.

Although both problems retrieve data subject to constraints, they answer fundamentally different questions. Constrained retrieval asks \emph{``Which existing objects satisfy a set of constraints?''} Counterfactual querying asks \emph{``What feasible changes to the query instance would lead to a different outcome?''} The goal is thus not to retrieve constrained objects, but to identify actionable, outcome-changing alternatives together with the changes needed to reach them.

This leads to four concrete distinctions. First, the opposite-outcome predicate and admissibility constraints jointly \emph{define} the feasible search space, rather than filter a similarity search. Second, candidates are ranked by intervention cost $\omega$---the effort to change the query instance---which is generally different from the geometric distance used by retrieval indexes. Third, the answer is a set of complementary counterfactuals that jointly optimize a user-specified objective $\Omega$ (e.g., diversity, robustness, or fairness), rather than independently ranked tuples; our new experiment shows this is not reducible to top-$k$ ranking (Table~\ref{tab:quality}). Fourth, each returned counterfactual carries an interpretable edit $\Delta(\mathbf{x}_q,\mathbf{x})$ describing how the query instance should change, making the result actionable rather than merely retrieved.

To make this clearer, we revised Section~1 and added Table~\ref{tab:positioning}, which compares constrained retrieval, vector search, nearest-neighbor search, subset selection, FACE~\cite{poyiadzi2020face}, and DcCA in terms of their query semantics rather than their implementation capabilities.}

\noindent\emph{Revision to Section~1:}
\begin{quote}
{\color{blue}Our goal is not to develop another counterfactual generation algorithm, but to establish \emph{counterfactual querying} as a first-class data management primitive. The challenge is therefore not generating a counterfactual, but designing the query abstraction, indexing mechanisms, optimization strategies, and execution model for counterfactual reasoning as a recurring database operation.

From this perspective, counterfactual querying differs fundamentally from existing retrieval pipelines. While existing systems compose vector search, filtering, reranking, and subset selection to retrieve similar objects, counterfactual querying asks \emph{which feasible interventions would change the outcome?} Accordingly, the opposite-outcome predicate and admissibility define the feasible search space, intervention cost $\omega$ replaces geometric similarity as the optimization objective, and the answer is a set of complementary counterfactuals that jointly optimizes a user-specified objective $\Omega$.

These semantics cannot be captured by composing existing retrieval operators. Treating retrieval, feasibility, intervention cost, and set selection as separate stages prevents holistic optimization and obscures the query semantics. We therefore elevate the opposite-outcome predicate, intervention cost $\omega$, and set objective $\Omega$ to first-class query constructs, enabling future data systems to optimize counterfactual reasoning as a native database operation.}
\end{quote}
